\documentclass[compsoc, a4paper]{IEEEtran}

\usepackage{fontspec}
\usepackage[catalan, bidi=basic, provide=*]{babel}

\babelprovide[import, onchar=ids fonts]{catalan}
\babelprovide[import, onchar=ids fonts]{english}

% Tipografia acadèmica
\babelfont{rm}{Noto Serif}

\usepackage{amsmath}
\usepackage{graphicx}
\usepackage{booktabs}
\usepackage{hyperref}
\usepackage{listings}

\linespread{1.2}

\begin{document}

\title{Resolució de Jocs Combinatoris Mitjançant Programació Dinàmica: El Joc de la Calculadora del 31 sota el Patró MVC}

\author{Josep~Oliver~Vallespir, ~Hugo~Valls~Sabater%
\IEEEcompsocitemizethanks{\IEEEcompsocthanksitem Estudiants del Grau en Enginyeria Informàtica a la Universitat de les Illes Balears (UIB).\protect\\
Assignatura: Algorismes Avançats.}%
}

\IEEEtitleabstractindextext{%
\begin{abstract}
Aquest document presenta el disseny, la formalització matemàtica i la implementació del joc de la calculadora del 31 sota un enfocament de resolució dinàmica amb memoització. S'hi detalla l'aplicació d'un algorisme de programació dinàmica generalitzat per a $N$ jugadors, així com la unificació d'estructures de dades primitives en objectes de domini (\texttt{EstatCalculat}) seguint principis de Programació Orientada a Objectes (POO). A més, es descriu l'arquitectura d'una interfície gràfica d'usuari interactiva altament desacoblada basada en el patró Model-Vista-Controlador (MVC), amb suport per a la col·locació aleatòria de tecles, detecció primerenca de derrota inevitable (\textit{Zugzwang}) i un sistema de control de fils concurrent per a l'execució de prediccions en temps real sense bloqueig de la interfície.
\end{abstract}

\begin{IEEEkeywords}
Programació Dinàmica, Jocs Combinatoris, Memoització, Inducció cap enrere, Disseny Orientat a Objectes, MVC, Concurrència, Zugzwang, Teclat dinàmic, Teoria de Jocs.
\end{IEEEkeywords}
}

\maketitle
\IEEEdisplaynontitleabstractindextext
\IEEEpeerreviewmaketitle

% ---------------------------------------------------
% SECCIÓ 1: INTRODUCCIÓ
% ---------------------------------------------------
\section{Introducció}

En el camp de la informàtica teòrica i la intel·ligència artificial, la resolució automatitzada de jocs combinatoris finits amb informació perfecta constitueix un dels pilars de la Teoria de Jocs. El joc analitzat en aquesta pràctica, conegut com a \textit{Joc de la Calculadora del 31}, és un joc d'oposició directa on dos o més jugadors prenen la decisió de quin nombre sumar a un acumulador central que té la calculadora. El jugador que arriba o supera un límit determinat (per defecte 31) perd de forma immediata la partida. La particularitat algorísmica i el que fa atractiu aquest joc és la presència de restriccions de moviment espacials: dels nombres del teclat (normalment disposats de l'1 al 9 en format de calculadora estàndard), el jugador en el seu torn només pot elegir-ne quatre, determinats per la jugada de l'oponent anterior. En concret, el jugador només pot escollir aquells nombres que es troben en la mateixa fila o columna física del teclat que el darrer nombre escollit, amb l'excepció de no poder repetir el mateix nombre.

El nucli del problema rau a determinar, donats uns paràmetres inicials de joc (suma actual, darrer nombre premut, límit establert i nombre de jugadors actius), quin dels jugadors acabarà perdent la partida suposant que tots juguen amb racionalitat perfecta (és a dir, aplicant una estratègia òptima). Aquest escenari es presta de forma immillorable a l'aplicació de la \textbf{Programació Dinàmica (PD)}, atès que el joc presenta una subestructura òptima clara i solapament de subproblemes. Per evitar la repetició ineficient de càlculs exponencials, s'utilitza una taula de memoització que emmagatzema les decisions i els camins resolts.

Més enllà de la pura teoria algorísmica, aquesta entrega aborda reptes crucials de disseny de programari. Tradicionalment, les estructures de memoització s'implementaven mitjançant matrius paral·leles disperses per a cada atribut a registrar. En aquesta pràctica, s'ha aplicat un enfocament orientat a objectes, encapsulant la informació de l'estat en una classe immutable. A més, s'ha dissenyat una interfície gràfica interactiva seguint una separació estricta en el patró \textbf{Model-Vista-Controlador (MVC)} absolutament pur, garantint la independència de la presentació respecte de la lògica de negoci i una concurrència eficient que executa la intel·ligència de joc en segon pla per evitar bloquejos visuals.

% ---------------------------------------------------
% SECCIÓ 2: ARQUITECTURA DE PROGRAMA I PATRO MVC
% ---------------------------------------------------
\section{Arquitectura del Programa i Patró MVC}

\subsection{Separació de Responsabilitats i Desacoblament pur}

L'arquitectura del programari s'ha estructurat seguint rigorosament el patró de disseny \textbf{Model-Vista-Controlador (MVC)}, assegurant un nivell màxim de cohesió interna i un acoblament pràcticament nul entre la visualització i la lògica operacional de la calculadora:

\begin{itemize}
    \item \textbf{Model (\texttt{model}):} El model encapsula les regles matemàtiques, l'estat del teclat i l'algorisme de programació dinàmica. Està conformat per les classes \texttt{Joc}, \texttt{Teclat} i \texttt{EstatCalculat}. Un aspecte fonamental és que cap d'aquestes classes conté cap referència a components gràfics (com Swing) ni coneix la presència de la interfície d'usuari. Això en facilita la reutilització en altres contextos (com línia de comandes o aplicacions web) i la realització de tests unitaris automatitzats de forma aïllada.
    \item \textbf{Vista (\texttt{vista}):} Està representada per la classe \texttt{Vista}. La seva única responsabilitat és dibuixar els components visuals i exposar controls a l'usuari. En la versió final d'aquest projecte, s'ha assolit un \textbf{desacoblament de llibre del 100\%}: la vista és completament agnòstica respecte a les classes del Model. La signatura del mètode de dibuixat del tauler de botons és:
    \begin{lstlisting}[language=Java, basicstyle=\tiny\ttfamily]
public void actualitzarTeclat(int[][] grid, int w, int h, int maxNombre)
    \end{lstlisting}
    D'aquesta manera, la Vista rep únicament matrius de tipus primitiu d'enters (\texttt{int[][]}) i desconeix per complet la classe de domini \texttt{Teclat}. Les operacions d'instanciació o càrrega d'estructures estan totalment vetades a la vista.
    \item \textbf{Controlador (\texttt{controlador}):} Representat per la classe \texttt{Controlador}, és el cervell connector de l'aplicació. Intercepta els clics als botons de la vista, actualitza les dimensions, re-genera el teclat (incloent-hi els barrejats aleatoris), executa l'algorisme del model i transfereix els resultats de tornada a les etiquetes de text de la vista de forma asíncrona.
\end{itemize}

\subsection{El Cicle de Comunicació Gràfica}

La comunicació entre les capes s'estableix de forma unidireccional i guiada per esdeveniments, evitant l'acoblament cíclic. El Controlador és qui s'encarrega de subscriure's a les accions de la vista mitjançant mètodes de registre de listeners gràfics. Durant l'usuari interactua amb un combobox (com la mida de la graella), la Vista notifica el Controlador mitjançant el mètode \texttt{setControladorDimensions}, i és el propi Controlador qui recrea el teclat en el Model, extreu la informació de coordenades i la envia a la Vista perquè redibuixi els botons interactius físics.

\subsection{Concurrència i Multi-threading de la IA}

Atès que l'algorisme de programació dinàmica per a jocs amb grans límits de puntuació o taulers molt amplis pot requerir milers de crides recursives, el càlcul de predicció del guanyador podria bloquejar el fil de processament visual principal de Java (\textit{Event Dispatch Thread} o EDT), congelant la pantalla i fent que l'aplicació sembli inestable. 

Per resoldre-ho, s'ha implementat concurrència:
\begin{itemize}
    \item \textbf{Fil de fons (Background Thread):} Cada vegada que es requereix preveure el guanyador o calcular l'entropia del joc, el Controlador engega un fil de treball aïllat (\texttt{new Thread}).
    \item \textbf{Interacció asíncrona:} Durant la computació, la vista és informada d'un estat ocupat (\texttt{setProcessant(true)}), mostrant un cursor de càrrega i deshabilitant els botons de configuració per evitar estats inconsistents.
    \item \textbf{SwingUtilities:} Un cop la programació dinàmica troba el resultat, aquest es transfereix de nou a l'EDT mitjançant crides a \texttt{SwingUtilities.invokeLater()} per actualitzar de manera segura els textos visuals i encendre el teclat, garantint que el programa sigui altament reactiu.
\end{itemize}

% ---------------------------------------------------
% SECCIÓ 3: DISSENY DE LA INTERFÍCIE D'USUARI (VISTA)
% ---------------------------------------------------
\section{Disseny de la Interfície d'Usuari (Vista)}

La interfície gràfica (GUI) ha estat desenvolupada utilitzant la biblioteca \textit{Swing} de Java, seguint els principis de disseny centrat en l'usuari. L'objectiu ha estat crear un entorn de treball net, professional i que permeti monitorar en tot moment l'estat del procés de joc. La pantalla principal es divideix en tres grans àrees funcionals:

\subsection{Panell Lateral d'Ajustaments i Controls}
Aquest panell reuneix tots els selectors que configuren l'inici de la simulació de la partida, organitzats de forma vertical en el lateral de la finestra:
\begin{itemize}
    \item \textbf{Selectors Numèrics (JSpinners):} Permeten fixar la \textit{Suma Inicial} (amb la qual comença el recompte central de la calculadora), el \textit{Límit} de la partida (que per defecte és 31, superat el qual es perd la partida), el \textit{Darrer Nombre} premut (que actua com la llavor dels moviments vàlids inicials), i el \textit{Nombre de Jugadors} ($N \ge 2$) de la partida.
    \item \textbf{Configuració del Tauler (JComboBox):} Selector dinàmic de mides de la graella de la calculadora. Permet passar a l'acte de disposicions de $3 \times 3$ (nombres de l'1 al 9), $3 \times 4$ o $4 \times 4$.
    \item \textbf{Targeta d'Accions:} Botons clau de simulació: "Calcular Guanyador" per llançar la IA en mode no interactiu, "Restablir Valors" per netejar la taula de joc, i l'acció especial **"Mesclar Teclat (Aleatori)"** que desordena espacialment els nombres.
\end{itemize}

\subsection{Panell Central de la Calculadora (Teclat Interactiu)}
L'espai central representa el teclat numèric interactiu. Es genera dinàmicament utilitzant un \texttt{GridLayout} les dimensions del qual s'ajusten immediatament a les files i columnes seleccionades. 
Durant les partides gràfiques interactives:
- Els botons corresponents a moviments vàlids (mateixa fila o columna que el darrer premut) es mantenen habilitats amb un fons gris clar.
- Els botons no vàlids es deshabiliten de forma immediata, agafant un color gris fosc i impedint el seu clic.
- La tecla seleccionada en el darrer moviment es pinta amb un ressaltat groc per permetre a l'usuari identificar amb claredat el focus de la jugada anterior.

\subsection{Panell Inferior i Recompte d'Estat}
Conté barres gràfiques de progrés i rètols informatius per oferir informació en temps real:
- **Estat de la partida:** Indica si s'està processant la predició de la IA, si s'ha fet un canvi de graella o si la partida ha finalitzat.
- **Predició de la IA:** Mostra en color verd brillant quin és el jugador destinat a perdre segons la programació dinàmica, oferint a l'usuari informació en directe de la "partida perfecta".
- **Comptador de camins:** Indica el nombre total de camins o partides possibles diferents que porten a una decisió des d'aquest estat, derivat directament del recompte de combinacions de la DP.

% ---------------------------------------------------
% SECCIÓ 4: IMPLEMENTACIÓ TÈCNICA DEL MODEL
% ---------------------------------------------------
\section{Implementació Tècnica del Model}

\subsection{L'Algorisme de Programació Dinàmica i Recursivitat}

El motor del model és el càlcul del guanyador mitjançant Programació Dinàmica. Donada la naturalesa seqüencial del joc, l'estat es pot identificar completament amb dues variables en cada torn:
1. $S$: La suma acumulada en aquell moment ($0 \le S \le L$).
2. $D$: El darrer nombre escollit ($0 \le D \le M$), on $0$ representa l'estat inicial del joc (on no hi ha restriccions espacials).

Es defineix la funció recursiva del perdedor relatiu, $P(S, D)$, com l'índex del jugador (dins de la seqüència local de torns, on $0$ és el jugador que té el torn actiu en aquest moment) que està destinat a perdre. La formulació matemàtica segueix el principi de la inducció cap enrere:

- **Cas base (Límit superat):** Si la suma actual és superior o igual al límit de derrota ($S \ge L$), el jugador que acaba de jugar i ha superat el límit perd de forma immediata. Com que la derrota s'avalua al principi del següent torn, el jugador que ara té el torn és qui rebrà la derrota immediata sense poder fer cap moviment:
  $$P(S, D) = 0 \quad \text{si } S \ge L$$

- **Pas recursiu (Tria de moviments òptims):** Si $S < L$, el jugador actiu ha de triar un moviment vàlid $v$ del conjunt de moviments permesos pel teclat a partir de $D$, anomenat $V(D)$. Cada elecció d'una tecla $v$ augmenta la suma a $S + v$ i cedeix el torn al següent jugador.
  Per a cada moviment $v \in V(D)$, es calcula recursivament qui perd en el següent estat, $P(S+v, v)$. Aquest perdedor es calcula des de la perspectiva del jugador que anirà després en el torn. Per tant, per traduir aquest perdedor relatiu de l'estat següent a la perspectiva del jugador actiu en aquest estat, cal aplicar un desplaçament circular en funció de les $N$ posicions de la taula de jugadors:
  $$P_{\text{candidat}}(v) = (P(S+v, v) + 1) \bmod N$$
  
  Com que cada jugador actua com a mínim amb racionalitat absoluta per evitar la seva pròpia derrota, el jugador actiu triarà aquell moviment $v$ que maximitzi les seves opcions de supervivència:
  1. Si hi ha com a mínim un moviment $v$ tal que $P_{\text{candidat}}(v) \neq 0$ (és a dir, triant $v$, el perdedor és un altre jugador i no ell), escollirà el moviment que allunyi al màxim la derrota d'ell (habitualment intentant que perdi el jugador següent en la seqüència, és a dir, $P_{\text{candidat}}(v) = 1$, per exercir màxima pressió).
  2. Si per desgràcia tots els moviments possibles $v \in V(D)$ porten a $P_{\text{candidat}}(v) = 0$ (és a dir, no importa què jugui, el perdedor serà irremeiablement ell mateix), llavors el jugador està en una posició de derrota inevitable (\textit{Zugzwang}), i qualsevol elecció que faci el farà perdre. En aquest cas, $P(S, D) = 0$.

De forma simultània, per calcular el nombre de partides o camins únics possibles que es poden jugar a partir d'un estat, es defineix la funció $C(S, D)$:
- $C(S, D) = 1$ si $S \ge L$ (una única seqüència terminal).
- $C(S, D) = \sum_{v \in V(D)} C(S+v, v)$ si $S < L$.

\subsection{Disseny POO i Memoització: La classe EstatCalculat}

En versions inicials de l'algorisme, la memoització es realitzava mitjançant la definició de tres matrius bidimensionals disperses independents de tipus primitius: `boolean[][] calculat`, `int[][] dpPerdedor` i `long[][] dpPartides`. Aquest patró és conegut com a codi dispers (*spaghetti*), ja que no respecta l'encapsulació i fa molt difícil el manteniment del codi o l'ampliació de nous atributs d'estat de la partida.

Per resoldre aquest problema, s'ha aplicat un disseny de Programació Orientada a Objectes sòlid. Hem creat la classe immutable \texttt{EstatCalculat}:

\begin{lstlisting}[language=Java, basicstyle=\tiny\ttfamily]
package model;

public class EstatCalculat {
    private final int perdedor;
    private final long partides;

    public EstatCalculat(int perdedor, long partides) {
        this.perdedor = perdedor;
        this.partides = partides;
    }

    public int getPerdedor() { return perdedor; }
    public long getPartides() { return partides; }
}
\end{lstlisting}

Així, la taula de memoització dins de \texttt{Joc.java} es redueix a una única graella de referències d'objecte:
\begin{lstlisting}[language=Java, basicstyle=\tiny\ttfamily]
private final EstatCalculat[][] dp;
\end{lstlisting}

Un estat particular $(S, D)$ es considera **no calculat** si `dp[S][D] == null`. Això elimina la necessitat d'una matriu booleana auxiliar per separat. Quan l'algorisme calcula el resultat per primera vegada, s'instancia un objecte de forma síncrona i es guarda a la taula. Aquest disseny redueix l'acoblament, facilita la recollida d'escombraries de la JVM gràcies a l'ús d'objectes locals de curta durada, i respecta els pilars fonamentals de la POO.

% ---------------------------------------------------
% SECCIÓ 5: GESTIÓ DINÀMICA DE LA GRAELLA (TECLAT)
% ---------------------------------------------------
\section{Gestió Dinàmica de la Graella Física (Teclat)}

\subsection{La Disposició Estàndard de Calculadora}

Per defecte, les calculadores físiques utilitzen una graella numèrica creixent de baix a dalt i d'esquerra a dreta. Hem implementat aquesta disposició a \texttt{Teclat.java}. Si denotem $W$ com l'amplada (columnes) i $H$ com l'alçada (files), el valor d'una tecla en la posició de fila $r$ i columna $c$ (indexades des de $0$) es calcula de forma espacial mitjançant la fórmula:
$$V(r, c) = (H - 1 - r) \times W + c + 1$$

Això genera exactament el disseny estàndard per a mides de $3 \times 3$:
```
7  8  9
4  5  6
1  2  3
```

\subsection{Determinació Espacial de Moviments vàlids en $O(1)$}

La principal restricció espacial del joc imposa que el següent moviment s'ha de triar exclusivament de la mateixa fila o columna física on es trobava el nombre premut anteriorment. 

Per tal d'avaluar aquesta restricció en temps constant $O(1)$ sinó haver de recórrer tota la matriu en cada jugada, la classe \texttt{Teclat} pre-calcula i emmagatzema les coordenades físiques de cada tecla en el moment de la seva construcció. S'utilitzen dues estructures de dades d'indexació immediata (vectors directes indexats pel valor del propi nombre menys u):
- `int[] filaDe`: Desa la fila física de cada valor.
- `int[] colDe`: Desa la columna física de cada valor.

Així, donat un nombre $D$, els moviments vàlids es poden extreure de manera extremadament ràpida avaluant directament si la distància de fila o columna és zero:
\begin{lstlisting}[language=Java, basicstyle=\tiny\ttfamily]
public ArrayList<Integer> getMovimentsValids(int darrerNombre) {
    ArrayList<Integer> valids = new ArrayList<>();
    if (darrerNombre == 0) {
        // Estat inicial: es poden premer totes les tecles
        for (int i = 1; i <= maxNombre; i++) valids.add(i);
        return valids;
    }
    int r0 = filaDe[darrerNombre - 1];
    int c0 = colDe[darrerNombre - 1];
    for (int v = 1; v <= maxNombre; v++) {
        if (v == darrerNombre) continue; // No es pot repetir el mateix
        int rv = filaDe[v - 1];
        int cv = colDe[v - 1];
        if (rv == r0 || cv == c0) {
            valids.add(v);
        }
    }
    return valids;
}
\end{lstlisting}

\subsection{Barrejat Aleatori de la Graella}

Per donar compliment al requisit d'innovació o variació de partida, s'ha implementat l'opció de **col·locació aleatòria de nombres** a la graella. 
Quan l'usuari activa aquesta opció (mitjançant el botó lateral), s'aplica una permutació aleatòria de la llista de nombres del $1$ al $maxNombre$ (utilitzant l'algorisme de barrejat de Fisher-Yates) i es col·loquen de forma desordenada en les posicions físiques de la matriu del teclat. 

El càlcul de moviments vàlids és completament dinàmic: com que s'actualitzen les coordenades `filaDe` i `colDe` segons la nova disposició barrejada, la restricció de la mateixa fila i columna gràfica es calcula perfectament sobre la base dels nous veïns espacials desordenats. Això altera completament els camins del joc de la calculadora, fent que cada combinació barrejada sigui un repte matemàtic diferent i molt més divertit.

% ---------------------------------------------------
% SECCIÓ 6: GENERALITZACIÓ A N JUGADORS I LÒGICA
% ---------------------------------------------------
\section{Generalització a N Jugadors i Lògica de Joc}

\subsection{Gestió de Torns Circulars per a N Jugadors}

L'algorisme està completament generalitzat per a un nombre arbitrari de jugadors ($N \ge 2$), trencant la limitació històrica de partides exclusives per a dues persones. El torn transcorre de manera seqüencial i circular: del jugador $k$, el torn passa al jugador $(k \bmod N) + 1$. 

La generalització es reflecteix en la inducció del perdedor de la DP. Quan un jugador fa un moviment que cedeix un estat on el perdedor és $P_{\text{seg}}$ (de tipus relatiu, on $0$ és el jugador següent), per al jugador actual aquest perdedor és desplaçat circularment:
$$P_{\text{actual}} = (P_{\text{seg}} + 1) \bmod N$$
Això assegura que per a qualsevol nombre de jugadors actius, es prevegi de forma exacte quin és el jugador absolut que acabarà superant el límit i perdent la partida.

\subsection{Detecció Primerenca de Defectes Inevitables (Zugzwang)}

En el disseny original del joc, el torn passava d'un jugador a un l'altre de manera indefinida fins que algú premia un botó que feia superar el límit establert, moment en el qual perdia. Això obligava a fer moviments de forma manual que eren destructius de forma òbvia.

Per millorar significativament la interactivitat del joc, hem implementat un control de **Defecte Inevitable (\textit{Zugzwang})** al \texttt{Controlador.java}. Abans de cedir el torn al següent jugador, el programa avalua les opcions espacials del següent jugador:
- Si **tots** els moviments de fila i columna física permesos fan que la nova suma superi o iguali el límit establert, és impossible que el jugador següent faci cap moviment segur.
- En detectar aquest estat de bloqueig, la partida **s'atura immediatament de forma automàtica**: es deshabilita el teclat per complet i s'actualitza l'etiqueta de resultats: *"Jugador X no té moviments segurs! HA PERDUT"*. 
Això evita jugades absurdes i dona retroalimentació de derrota primerenca instantània.

\subsection{Reinici Robust sense Càlculs Residuals}

Un repte especial en la programació visual d'Swing ha estat el botó de **"Restablir Valors"**. En prémer aquest botó, s'inicialitzen tots els controls de la pantalla als seus valors per defecte (graella a $3 \times 3$, límit a 31, suma a 0). Com que el canvi del combobox llança esdeveniments asíncrons a la cua de Swing, es produïa una carrera d'esdeveniments que calculava de forma asíncrona la solució enmig de la neteja, deixant resultats residuals a la pantalla en lloc de deixar-la buida.

Per solucionar-ho, hem definit un flag lògic `ignorarEvents` al Controlador. En fer clic a restablir:
1. S'activa `ignorarEvents = true;`
2. Es netegen tots els controls de la vista. Swing encolarà els esdeveniments de canvi de valors.
3. Els esdeveniments són interceptats i el controlador retorna immediatament sense executar cap programació dinàmica, mantenint les etiquetes netes en l'estat *"Pendent de càlcul..."*.
4. S'utilitza `SwingUtilities.invokeLater(() -> ignorarEvents = false)` per restablir el flag només quan Swing hagi buidat per complet la cua de neteja, assegurant un reinici net, robust i d'alta qualitat.

% ---------------------------------------------------
% SECCIÓ 7: OPTIMALITAT I CORRECTESA DE L'ALGORISME
% ---------------------------------------------------
\section{Optimalitat i Correctesa de la Resolució}

\subsection{El Principi d'Optimalitat sota Perfecta Racionalitat}

La determinació del guanyador a través de la Programació Dinàmica es basa en el supòsit de perfecta racionalitat dels jugadors. En no haver-hi elements de sort (com daus), el joc del 31 es pot formalitzar com un graf dirigit acíclic d'estats de joc. 

L'aplicació d'inducció cap enrere garanteix la correctesa de la decisió. Si un estat té un perdedor relatiu de $0$ per a tots els moviments, l'estat és guanyador per al rival. L'algorisme és matemàticament exhaustiu i correcte perquè avalua absolutament totes les ramificacions de moviments permesos sobre la base de la subestructura espacial del teclat.

\subsection{Anàlisi de l'Efecte Zugzwang i Simetries}

L'efecte de *Zugzwang* es dóna en aquells estats on tenir el torn és un desavantatge absolut (qualsevol moviment vàlid porta a la derrota). L'algorisme de memoització detecta aquests estats a l'acte: quan un estat té el valor recursiu de perdedor relatiu igual a $0$, s'està davant d'una posició on el jugador actiu perdrà de forma inevitable si el contrincant juga correctament, posant de manifest la rigidesa espacial del disseny bidimensional del teclat.

% ---------------------------------------------------
% SECCIÓ 8: CAS PRÀCTIC: FUNCIONAMENT PAS A PAS
% ---------------------------------------------------
\section{Cas Pràctic: Funcionament Pas a Pas}

Per il·lustrar empíricament la correctesa i la lògica del nostre algorisme recursiu, analitzem un escenari simplificat i real de partida.

\subsection{Escenari Interactiu en Teclat Estàndard $3 \times 3$}

Partim del teclat estàndard de calculadora:
```
7  8  9
4  5  6
1  2  3
```

Configurem la partida amb els següents paràmetres inicials de prova:
- **Límit d'acumulació:** 31
- **Nombre de jugadors:** 2
- **Suma acumulada actual:** 28
- **Darrer nombre premut:** 6 (correspon al focus grog de la pantalla)

\subsection{Determinació de Moviments vàlids}
El darrer nombre premut és 6. La seva posició en la graella és la fila 1 i la columna 2. Els veïns espacials de la seva mateixa fila (4, 5) o columna (9, 3) (excloent el propi 6) són el conjunt de moviments permesos:
$$V(6) = \{3, 4, 5, 9\}$$

\subsection{Avaluació del Torn i Detecció de Zugzwang}
Suposem que és el torn del **Jugador 1**. Analitzem els resultats de cadascun dels moviments possibles de $V(6)$ que té sobre la taula:
1. **Si prem 3:** La nova suma és $28 + 3 = 31$. Com que arriba al límit, el Jugador 1 perd immediatament.
2. **Si prem 4:** La nova suma és $28 + 4 = 32$. Supera el límit, perd.
3. **Si prem 5:** La nova suma és $28 + 5 = 33$. Supera el límit, perd.
4. **Si prem 9:** La nova suma és $28 + 9 = 37$. Supera el límit, perd.

Com que **tots** els moviments vàlids disponibles des del 6 porten a una nova suma superior o igual al límit de 31, el Jugador 1 es troba en un estat de derrota inevitable de tipus *Zugzwang*. 

Gràcies a la nostra millora en la detecció primerenca del controlador, en lloc de permetre al Jugador 1 pitjar un botó destructiu per força, l'aplicació interactiva atura el joc a l'acte: deshabilita de color gris fosc el teclat, posa l'estat en color vermell com a *"Fi del joc (Derrota inevitable)"*, i mostra a la pantalla:
**"Jugador 1 no té moviments segurs! HA PERDUT"**.

Aquesta traça es veu reflectida de forma neta i simple en el terminal del sistema, confirmant que el comportament de la interfície gràfica i el nucli del model estan 100\% sincronitzats i són correctes.

% ---------------------------------------------------
% SECCIÓ 9: ANÀLISI DE COMPLEXITAT
% ---------------------------------------------------
\section{Anàlisi de Complexitat}

\subsection{Complexitat Temporal}

La complexitat temporal del nostre algorisme és excel·lent gràcies a l'aplicació de la Programació Dinàmica com a mètode d'optimització:

- **Sense memoització (Recursió simple):** El nombre d'estats creixeria de forma exponencial en funció del límit de la partida, representat com $O(4^L)$, fent impossible calcular partides amb límits superiors a 40.
- **Amb memoització (PD):** Cada estat definit pel parell $(S, D)$ es calcula una única vegada. Hi ha com a màxim $L$ sumes diferents ($0 \le S \le L$) i $M$ darrers nombres possibles ($0 \le D \le M$), on $M$ és la quantitat de tecles del teclat (9 per a $3 \times 3$, 16 per a $4 \times 4$). 
  Per a cada estat no calculat, l'algorisme realitza un bucle sobre els $M$ moviments vàlids del teclat. Per tant, la complexitat temporal de pitjor cas de la construcció de la taula DP és de tipus polinomial:
  $$T(L, M) = O(L \times M^2)$$

Per a un límit típic de $L=31$ i teclat de $3 \times 3$ ($M=9$), el nombre màxim d'estats a calcular és inferior a $31 \times 10 = 310$ estats, que requereixen menys d'un mil·lisegon de temps de CPU de la JVM. Això fa que el càlcul de la IA sigui instantani i extremadament eficient, fins i tot per a sumes límit molt elevades.

\subsection{Complexitat Espacial}

L'espai de memòria RAM utilitzat pel programa depèn directament de la taula de memoització. Com que hem consolidat la programació dinàmica en una sola graella d'objectes bidimensional `EstatCalculat[L][M]`, la complexitat espacial és strictly lineal respecte al límit de la suma i el nombre de botons del teclat:
$$S(L, M) = O(L \times M)$$

Atès que les cèl·les s'omplen de forma mandrosa (*lazy*) només a mesura que els estats són visitats de forma efectiva pel recorregut recursiu, el consum espacial real de la JVM en execució és extremadament reduït (uns pocs bytes de metadades per a les referències), cosa que garanteix un consum totalment estable de memòria.

% ---------------------------------------------------
% SECCIÓ 10: CONCLUSIÓ
% ---------------------------------------------------
\section{Conclusió}

La realització de la pràctica del Joc del 31 ha posat de manifest la importància cabdal d'aplicar patrons d'arquitectura nets i estructures de dades coherents en el desenvolupament de programari complex. En lloc de limitar-nos a resoldre el problema lògic a nivell matemàtic, hem anat un pas més enllà per garantir la màxima qualitat d'enginyeria.

En primer lloc, hem integrat conceptes avançats de la Programació Orientada a Objectes (POO) a la Programació Dinàmica. La creació de la classe immutable \texttt{EstatCalculat} ens ha permès unificar les dades lògiques disperses de la memoització, eliminant variables globals i matrius primitives paral·leles redundants. Aquest disseny no només fa que el codi sigui molt més modular, llegible i fàcil de mantenir, sinó que optimitza l'ús de la memòria de forma molt intel·ligent.

En segon lloc, s'ha assolit un desacoblament de dependències absolut en el patró MVC. En fer que la vista rebi exclusivament valors i graelles primitives per paràmetre, hem garantit que la representació física i gràfica sigui completament autònoma del domini dels càlculs. Això s'ha complementat de forma excel·lent amb una interfície gràfice concurrent i segura en Swing, que executa la IA en segon pla per evitar bloquejos i utilitza fils independents i flags com `ignorarEvents` per resoldre amb èxit carreres crítiques d'esdeveniments.

Finalment, la incorporació de variacions de joc com el teclat dinàmic barrejat de forma espacial i el control del *Zugzwang* demostra com els algorismes asimptòtics i la programació dinàmica es poden integrar en entorns interactius reals per oferir una experiència dinàmica, ràpida i de gran qualitat de programari, complint amb escreix els estàndards exigits per l'assignatura d'Algorismes Avançats.

% ---------------------------------------------------
% BIBLIOGRAFIA
% ---------------------------------------------------
\begin{thebibliography}{1}

\bibitem{aa}
Assignatura d'Algorismes Avançats, Universitat de les Illes Balears, 2026.

\bibitem{gemini}
\textit{Gemini}, model 3.5 Pro, Google.

\bibitem{dynamic}
Bellman, R., ``Dynamic Programming,'' \textit{Science}, vol. 153, no. 3731, pp. 34-37, 1966.

\bibitem{swing_concurrency}
Oracle Corporation, ``Concurrency in Swing,'' \textit{The Java Tutorials}. [En línia]. Disponible a: \url{https://docs.oracle.com/javase/tutorial/uiswing/concurrency/index.html}.

\end{thebibliography}

\end{document}
